\section{Executed pre-study validation}
\label{app:diagnostics}
The following are measured implementation and feasibility diagnostics, separate from the completed confirmatory matrix. They establish neither method rankings nor fitted-composition population coverage.

\paragraph{Calibration variability.}
With $n=499$, $\alpha=.1$, $k=450$ and 2,000 test draws, the saved 10,000-repeat continuous-score simulation gives mean population coverage $.9001285$ and variance $.00017523$, compared with the analytic mean $.9$ and variance $.00017964$ in Equation~\eqref{eq:beta}. Empirical finite-test coverage variance is $.00022167$ versus $.00022455$ from Equation~\eqref{eq:beta-binomial}. Population coverage conditional on a realized calibration set falls below $.9$ in $.4835$ of these draws. This illustrates calibration variability rather than a marginal-guarantee violation. Figure~\ref{fig:calibration-diagnostic} reports the actual empirical CDFs with their separate Monte Carlo uncertainty bands. Its inputs are saved scores/count-derived coverage samples, not fitted-composition study trials.

\paragraph{Area integration and pilot computation.}
Figure~\ref{fig:integration-diagnostic} covers 66 fixed TV/L2/Hellinger/ILR cases, including smoothed zero-arm scores. This supersedes the earlier 51-case diagnostic while preserving its historical record. At $2^{14}$ points, the largest discrepancy of the four-scramble mean from a separate $2^{21}$-point direct uniform-simplex reference is $.000856$; at $2^{16}$ it is $.000115$. The reference also has finite numerical error, and these maxima are not universal error bounds. A separate paired high-resolution audit on saved pilot fits computes 304 matched-location contrasts under Equation~\eqref{eq:allowance}; its largest empirical allowance is $.002515$. The pilot's two-trial intervals are not interpreted as scientific comparisons.

The final-power pilot executes five $K=3$ mechanisms with two trials each, TV included, $2^{14}$ points and 32 area locations per group. It takes 14.105111 seconds and peaks at 180,316 KiB RSS in the supplied CPU environment. This is measured pilot cost, not the final study's cost. The source-pinned pilot metadata and saved rows/checkpoints are in \texttt{analysis/results/final\_candidate\_pilot/}; diagnostic and paired-audit records are under \texttt{final\_diagnostics/} and \texttt{paired\_gate/}. The analyst's approximately 47-minute main-run forecast, with a two-hour shared-host allowance, was a planning extrapolation; Appendix~\ref{sec:process} reports actual completed-study costs.

\paragraph{Checks, failures and changes.}
Thirteen implementation checks pass in the integrated checkout, including quantile/infinity branches, ILR isometry, true-zero original-target membership, absent-test-group aggregation, and independent interior/near-boundary volume checks. A missing $\sqrt K$ normalization in a working ILR Jacobian was caught and corrected before pilot computation. Early additive integration seeds overlapped between trial/location indices; the later source-pinned pilot uses separated keys. Exploratory pilot outputs are retained without relabeling them as final-code evidence. One audit launcher failed before computation because a sandbox temporary path was unavailable; the corrected job used the canonical shared config. These changes and failure records are supplied in \texttt{analysis/}; no confirmatory outcome was produced by the earlier faulty implementations.

\input{figures/diagnostics}
